\documentclass[11pt]{article}
\usepackage[margin=1in]{geometry}
\usepackage{amsmath,amssymb,amsthm}
\usepackage[T1]{fontenc}
\usepackage[utf8]{inputenc}
\usepackage{lmodern}
\usepackage[hidelinks]{hyperref}

\newtheorem{proposition}{Proposition}
\newcommand{\cb}{\operatorname{cb}}
\newcommand{\ck}{\operatorname{ck}}
\newcommand{\co}{\operatorname{conv}}
\newcommand{\supp}{s}

\title{Literature-Implied Full Negative Answer:\\
Positive Pettis Integrability Need Not Be Strong}
\author{Run \texttt{fa\_banach\_001}, agent \texttt{agent\_lane\_08}}
\date{9 August 2026}

\begin{document}
\maketitle

\section*{Status}

\textbf{Status: literature-implied full negative answer} to Problem 4.7 and to
the second question in Problem 4.8 of arXiv:1812.00597.  The decisive example
is Example 1.12 of Musia{\l} (2011); the implication to strong Pettis
integrability below appears not to be stated there and is identified directly
in this packet.

\section*{Source questions}

Candeloro, Di Piazza, Musia{\l}, and Sambucini ask on PDF page 11 of
\cite{source}:
\begin{quote}
\textbf{Problem 4.7.} Is each positive Pettis integrable multifunction
McShane integrable?

\textbf{Problem 4.8, second question.} Is each positive Pettis integrable
multifunction strongly Pettis integrable?
\end{quote}
Here a Pettis-integrable multifunction $\Gamma\colon[0,1]\to\cb(X)$ is
\emph{strongly Pettis integrable} when its Pettis primitive
$M_\Gamma$ is countably additive in the Hausdorff metric, that is, when it is
an $h$-multimeasure.  The same paper's Theorem 4.2 (PDF pages 9--10) states
that a $\cb(X)$-valued multifunction is McShane integrable if and only if it is
both strongly Pettis and Henstock integrable \cite{source}.

\section*{Supporting example}

Example 1.12 of Musia{\l} \cite[p.~779; PDF p.~11]{musial} chooses a scalarly
integrable $f\colon[0,1]\to c_0$ such that
\[
  T_f\colon c_0^*=\ell_1\longrightarrow L_1[0,1],
  \qquad T_f(x^*)=\langle x^*,f(\cdot)\rangle,
\]
is not weakly compact, and defines
\[
  \Gamma(t)=\co\{0,f(t)\}.
\]
The example records that $\Gamma$ is $\ck(c_0)$-valued, scalarly integrable,
and Pettis integrable in $\cb(c_0)$.  Moreover
\begin{equation}\label{eq:support}
  \supp(x^*,\Gamma(t))
  =\max\{0,\langle x^*,f(t)\rangle\}
  =\langle x^*,f(t)\rangle^+\geq 0,
\end{equation}
so $\Gamma$ is positive in precisely the sense used in the source question.

\section*{Identification}

\begin{proposition}
The multifunction $\Gamma$ in Musia{\l}'s Example 1.12 is Pettis integrable
but not strongly Pettis integrable.  Consequently it is not McShane
integrable.
\end{proposition}

\begin{proof}
Let $M_\Gamma$ be the $\cb(c_0)$-valued Pettis primitive of $\Gamma$.  Assume
for contradiction that it is an $h$-multimeasure.  Use the R\aa dstr\"om
support-function embedding
\[
  j\colon\cb(c_0)\longrightarrow \ell_\infty(B_{\ell_1}),
  \qquad j(C)(x^*)=\supp(x^*,C).
\]
This is an isometry from Hausdorff distance to the supremum norm and converts
Minkowski addition into vector addition.  Hence
\[
  m(E):=j(M_\Gamma(E))
\]
is a norm-countably additive $\ell_\infty(B_{\ell_1})$-valued vector measure.
In particular, if $E_n\downarrow\varnothing$, continuity of a vector measure
from above gives
\begin{equation}\label{eq:continuity}
  \lVert m(E_n)\rVert_\infty\longrightarrow 0.
\end{equation}

The defining identity for the Pettis primitive and
\eqref{eq:support} give, for every measurable $E$ and $x^*\in B_{\ell_1}$,
\[
  m(E)(x^*)
   =\supp(x^*,M_\Gamma(E))
   =\int_E \langle x^*,f(t)\rangle^+\,dt.
\]
All these coordinates are nonnegative.  Therefore
\begin{align*}
  \sup_{x^*\in B_{\ell_1}}
  \int_E |\langle x^*,f(t)\rangle|\,dt
  &\leq
  \sup_{x^*\in B_{\ell_1}}
  \int_E\langle x^*,f(t)\rangle^+\,dt \\
  &\quad+
  \sup_{x^*\in B_{\ell_1}}
  \int_E\langle -x^*,f(t)\rangle^+\,dt \\
  &\leq 2\lVert m(E)\rVert_\infty.
\end{align*}
Taking $E=[0,1]$ proves $L_1$-boundedness of $T_f(B_{\ell_1})$; taking
$E=E_n$ and using \eqref{eq:continuity} proves its uniform countable
additivity.  Equivalently, $T_f(B_{\ell_1})$ is uniformly integrable.
The Dunford--Pettis weak-compactness criterion in $L_1$ now implies that
$T_f(B_{\ell_1})$ is relatively weakly compact.  Thus $T_f$ is weakly compact,
contrary to the defining property in Example 1.12.  This contradiction proves
that $M_\Gamma$ is not an $h$-multimeasure.

Finally, Theorem 4.2 of \cite{source} implies that every McShane-integrable
$\cb(X)$-valued multifunction is strongly Pettis integrable.  Hence this
$\Gamma$ cannot be McShane integrable.
\end{proof}

\section*{Independent confirmation of the McShane answer}

Example 2.10 of Candeloro, Di Piazza, Musia{\l}, and Sambucini
\cite[PDF p.~9]{companion} independently starts with a Pettis-integrable but
non-McShane-integrable vector function $g$ and sets
$G(t)=\co\{0,g(t)\}$.  It explicitly concludes that $G$ is a
$\ck(X)$-valued Pettis-integrable multifunction containing $0$ everywhere,
hence positive, but is neither Henstock nor McShane integrable.  Thus the
negative answer to Problem 4.7 is also present in the later literature without
using the strong-Pettis implication above.

\section*{Provenance and scope}

The construction and the non-weak-compactness of $T_f$ are published in
Musia{\l}'s 2011 Example 1.12.  The supporting paper predates the source
questions and does not state the strong-Pettis consequence.  This packet is
therefore classified as an \emph{agent-identified literature implication},
not as a new counterexample produced by the run.

The result fully answers Problem 4.7 and the second question of Problem 4.8.
It does \emph{not} answer Problem 4.8's first question about positive
Henstock-integrable multifunctions possessing Henstock-integrable selections.
Nor does it answer the stronger question, posed as Question 2.14 of
\cite{companion}, asking for a positive Henstock- and Pettis-integrable
$\cb(c_0)$-valued multifunction that is not strongly Pettis integrable.
Nothing cited here shows that Musia{\l}'s Example 1.12 is Henstock integrable.

\section*{Review focus}

The only non-verbatim bridge needing expert verification is the passage from
Hausdorff countable additivity of $M_\Gamma$ to uniform integrability of
$T_f(B_{\ell_1})$.  The displayed calculation isolates that bridge: the
R\aa dstr\"om image is a norm-countably additive vector measure, and positivity
turns its norm continuity into uniform control of both positive and negative
parts of every scalar evaluation.

\begin{thebibliography}{9}

\bibitem{source}
D. Candeloro, L. Di Piazza, K. Musia{\l}, and A. R. Sambucini,
\emph{Integration of multifunctions with closed convex values in arbitrary
Banach spaces}, arXiv:1812.00597v4 (2019),
\url{https://arxiv.org/abs/1812.00597}.

\bibitem{musial}
K. Musia{\l},
\emph{Pettis Integrability of Multifunctions with Values in Arbitrary Banach
Spaces}, Journal of Convex Analysis 18 (2011), no.~3, 769--810,
\url{https://www.heldermann-verlag.de/jca/jca18/jca1004_b.pdf}.

\bibitem{companion}
D. Candeloro, L. Di Piazza, K. Musia{\l}, and A. R. Sambucini,
\emph{Multifunctions determined by integrable functions}, arXiv:1906.07019v2
(2020), \url{https://arxiv.org/abs/1906.07019}.

\end{thebibliography}

\end{document}
